% Spaltengruppen der Ergebnis-CSV (Abschnitt 3.5.1)
% Benoetigt in der Praeambel: \usepackage{booktabs,tabularx}
\begin{table}[htbp]
  \centering
  \small
  \caption{Spaltengruppen der Ergebnisdatei \texttt{results\_<Datum>.csv}. Jede Zeile entspricht einer Iteration eines Experiments.}
  \label{tab:csv-spalten}
  \begin{tabularx}{\textwidth}{@{}>{\raggedright\arraybackslash}p{2.6cm} >{\raggedright\arraybackslash}p{5.2cm} >{\raggedright\arraybackslash}X@{}}
    \toprule
    \textbf{Gruppe} & \textbf{Spalten} & \textbf{Inhalt} \\
    \midrule
    Experiment\-parameter
      & \texttt{task\_id}, \texttt{algorithm}, \texttt{function}, \texttt{dimension}, \texttt{scenario}, \texttt{intensity}, \texttt{repetition}, \texttt{f\_opt},~\dots
      & Einstellungen des Experiments aus dem Experimentplan, in jeder Zeile wiederholt \\
    \addlinespace
    Iterations\-werte
      & \texttt{iteration}, \texttt{next\_x}, \texttt{next\_y}, \texttt{next\_f\_true}, \texttt{acq\_value}
      & Vorgeschlagener Punkt, beobachteter und wahrer Funktionswert, Wert der Akquisitionsfunktion \\
    \addlinespace
    Regret
      & \texttt{best\_observed\_so\_far}, \texttt{best\_true\_so\_far}, \texttt{simple\_regret}, \texttt{instantaneous\_regret}, \texttt{cumulative\_regret}
      & Bisher bester beobachteter und wahrer Wert, Regret-Größen relativ zu $f_{\mathrm{opt}}$ \\
    \addlinespace
    Empfehlungs\-punkte und Metriken
      & \texttt{x\_rec\_<R>}, \texttt{M1\_<R>}, \dots, \texttt{M6\_<R>} \newline
        mit $R \in \{\texttt{BO}, \texttt{POM}, \texttt{ROBUST}\}$
      & Empfohlener Punkt je Empfehlungsregel und dessen Robustheitsmetriken M1--M6 \\
    \addlinespace
    Modell\-zustand
      & \texttt{prior\_offset}, \texttt{n\_train}, \texttt{mu\_train}, \texttt{sigma\_train}, \texttt{hyperparameters}
      & Prior-Fehler sowie Posterior-Mittelwert, -Standardabweichung und Hyperparameter an den Trainingspunkten (JSON) \\
    \addlinespace
    Startdaten
      & \texttt{start\_X}, \texttt{start\_f\_obs}, \texttt{start\_f\_true}
      & Startdesign mit beobachteten und wahren Werten, nur in Iteration~1 gefüllt (JSON) \\
    \addlinespace
    Laufzeit und Status
      & \texttt{iter\_time\_s}, \texttt{status}, \texttt{error}
      & Rechenzeit der Iteration in Sekunden; \texttt{ok} oder \texttt{error} mit Fehlermeldung \\
    \bottomrule
  \end{tabularx}
\end{table}
